\section*{Bab \ref{ch:b3:quantitative-geomorphology} --- Geomorfologi Kuantitatif}

\begin{solution}{exo:b3:quantitative-geomorphology:1}
$a/\sqrt{\pi\kappa t} = 0.30 - 0.04 = 0.26$, jadi $\kappa t = (2/0.26)^2/\pi =
\qty{18.8}{m^2}$, dan $t = 18.8/(5\times10^{-3}) = \qty{3.8}{ka}$, dalam
faktor dua ke kedua arah, mengingat rentang $\kappa$.
\end{solution}

\begin{solution}{exo:b3:quantitative-geomorphology:2}
$E = \kappa\,|\nabla^2 z| = 4\times10^{-3}\times0.012 =
\qty{4.8e-5}{m/yr}$, \qty{0.048}{mm/yr}. Relief $E L^2/(2\kappa) =
0.012\times80^2/2 = \qty{38.4}{m}$.
\end{solution}

\begin{solution}{exo:b3:quantitative-geomorphology:3}
$E = 2\times10^{-6}\times(25\times10^6)^{0.5}\times0.01 = 2\times10^{-6}
\times5000\times0.01 = \qty{1.0e-4}{m/yr}$, \qty{0.1}{mm/yr}. Karena $E$ bersatuan
\unit{m/yr}, $A^m$ bersatuan $\unit{m}^{2m}$ dan $S$ tanpa satuan, $K$ bersatuan
$\unit{m}^{1-2m}\,\unit{yr^{-1}}$: \unit{yr^{-1}} untuk $m = 0.5$.
\end{solution}

\begin{solution}{exo:b3:quantitative-geomorphology:4}
(a) Tetap: lapisan keras membuat sungai lebih curam di tempat ia tersingkap, dan undakan
tetap melekat pada lapisan itu. (b) Terpaku pada sesar: selama sesar bergeser,
sungai menjadi lebih curam di tempat ia memotong sesar itu. (c) Bergerak: pengangkatan
pesisir menurunkan \omterm{def:b3:quantitative-geomorphology:knickpoint}{dasar erosi} setiap sungai sekaligus; titik-titik tekuk
sedang bergerak ke hulu dan, seperti diprediksi teorema gelombang kinematik untuk $n
= 1$, semuanya berada pada satu kontur.
\end{solution}

\begin{solution}{exo:b3:quantitative-geomorphology:5}
$\theta = m/n = 0.5$. $S = (U/K)^{1/2}A^{-1/2} = (10^{-4}/10^{-9})^{1/2}
\times(10^7)^{-1/2} = 316\times3.16\times10^{-4} = 0.10$. Melipatgandakan $U$
mengalikan $k_s = (U/K)^{1/n}$ dengan $\sqrt2 = 1.41$; untuk $n = 1$ dengan 2:
makin besar $n$, makin kecil respons kecuraman terhadap pengangkatan.
\end{solution}

\begin{solution}{exo:b3:quantitative-geomorphology:6}\begingroup\emergencystretch=3em\relax 
Dengan $x$ diukur ke hulu dari muara, $\mathrm{d}x = -\mathrm{d}d$, sehingga $\chi(d)
= \int_d^L (A_0/c\,d'^{h})^{\theta}\,\mathrm{d}d' = A_0^{\theta} c^{-\theta}
\int_d^L d'^{-h\theta}\,\mathrm{d}d' = A_0^{\theta}c^{-\theta}(L^p -
d^p)/p$. Di sini $p = 0.19$, $A_0^{\theta}c^{-\theta} = 10^{2.7}/2^{0.45}$,
dan $\chi = \qty{2935}{m}$.
\par\endgroup\end{solution}

\begin{solution}{exo:b3:quantitative-geomorphology:7}
$c = 2\times10^{-6}\times(10^8)^{0.5} = \qty{0.02}{m/yr}$, \qty{20}{m/ka};
\qty{5}{km} dengan kecepatan itu butuh \qty{250}{ka}. Luas daerah tangkapan, dan bersamanya
kelajuan, berkurang ke arah hulu, sehingga \omterm{def:b3:quantitative-geomorphology:knickpoint}{titik tekuk} melambat saat naik.
\end{solution}

\begin{solution}{exo:b3:quantitative-geomorphology:8}
$\bar e = 0.40\times600 = \qty{240}{m}$; puncak naik sebesar $0.80\times240 =
\qty{192}{m}$; ketinggian rata-rata turun sebesar $240 - 192 = \qty{48}{m}$.
\end{solution}

\begin{solution}{exo:b3:quantitative-geomorphology:9}
$h = h_0\ln(P_0/E)$: $0.5\ln 6 = \qty{0.90}{m}$ dan $0.5\ln1.5 =
\qty{0.20}{m}$. Pada \qty{0.4}{mm/yr} erosi melampaui \omterm{def:b3:advanced-geochronology:cosmogenic}{laju produksi} tercepat,
$P_0$: tanah terkikis habis dan lereng menjadi batuan telanjang, penurunannya
dibatasi oleh pelapukan.
\end{solution}

\begin{solution}{exo:b3:quantitative-geomorphology:10}
Pada \omterm{def:b3:quantitative-geomorphology:steady-profile}{keadaan tunak} $\mathrm{d}q/\mathrm{d}x = E$ dan $q = 0$ di punggungan, sehingga
$q = E x$. Maka $\kappa S = E x\,(1 - S^2/S_c^2)$, persamaan kuadrat yang akar positifnya
adalah $S = [-\kappa + \sqrt{\kappa^2 + 4E^2x^2/S_c^2}\,]/(2Ex/S_c^2)$. Pada
20, 50, 100 dan \qty{200}{m}: 0.36, 0.68, 0.89 dan 1.03, dibandingkan dengan 0.4, 1.0,
2.0 dan 4.0 untuk difusi linear ($S = Ex/\kappa$). Bukit memiliki puncak yang
membulat, lalu sisi-sisi yang gradiennya cenderung ke $S_c$: lereng hampir datar, jauh
lebih rendah reliefnya daripada parabola.
\end{solution}

\begin{solution}{exo:b3:quantitative-geomorphology:11}
Menurut teorema gelombang kinematik, perubahan mencapai titik dengan koordinat
$\chi$ setelah $\chi/(KA_0^m)$; jaringan sudah sepenuhnya merespons ketika telah
mencapai hulu alur terakhir, dengan $\chi$ terbesar. $K A_0^m = 10^{-6}\times
1000 = \qty{e-3}{m/yr}$: batang utama merespons dalam \qty{5.46}{Ma}, jaringan
dalam \qty{7.24}{Ma}, sesuai dengan orde satu juta tahun yang diharapkan.
\end{solution}

\begin{solution}{exo:b3:quantitative-geomorphology:12}
Di atas pertemuan, gradien \omterm{def:b3:quantitative-geomorphology:steady-profile}{keadaan tunak} anak sungai adalah $S = (U/2K)A^{-\theta}$,
sehingga $\mathrm{d}z/\mathrm{d}\chi = (U/2K)A_0^{-\theta}$, setengah dari batang utama.
\omterm{def:b3:quantitative-geomorphology:knickpoint}{Titik tekuk} mencapai pertemuan ketika batang utama berada di $\chi_j$; setelah
tambahan waktu $\Delta t$, batang utama telah bergeser ke $\chi_k$ dengan $\chi_k -
\chi_j = K A_0^m\,\Delta t$, sedangkan anak sungai, dengan kelajuan $2KA^m$,
telah bergeser ke $\chi$ dengan $\chi - \chi_j = 2KA_0^m\,\Delta t = 2(\chi_k -
\chi_j)$.
\end{solution}

\begin{solution}{pb:b3:quantitative-geomorphology:1}
\textbf{1.} $4.5\times12000^{1.7} = \qty{3.87e7}{m^2}$, \qty{38.7}{km^2};
$4.5\times300^{1.7} = \qty{7.3e4}{m^2}$, \qty{0.073}{km^2}.

\textbf{2.} $\theta \approx -\log(0.007855/0.06641)/\log(37.08/0.388) =
0.928/1.980 = 0.47$.

\textbf{3.} Ya: dengan sebaran sekitar \qty{10}{\percent} pada dua belas titik,
$0.466$ tidak berbeda secara signifikan dari 0.5.

\textbf{4.} $0.06641\times\sqrt{388\,000} = \qty{41.4}{m}$;
$0.007855\times\sqrt{3.708\times10^7} = \qty{47.8}{m}$. Rata-rata geometrik,
\qty{45.1}{m}, memberikan taksiran $k_s$ untuk $\theta = 0.5$.

\textbf{5.} Gradien adalah selisih dua ketinggian, masing-masing tidak pasti
sebesar desimeter hingga meter, dibagi jarak yang pendek; pada ruas landai
galat merupakan bagian besar dari selisih itu. $\chi$ mengintegrasikan, bukan
mendiferensialkan.

\textbf{6.} $p = 1 - 1.7\times0.5 = 0.15$; $\chi = (1000/\sqrt{4.5})\,
(12000^{0.15} - 9000^{0.15})/0.15 = \qty{543.1}{m}$.

\textbf{7.} $(193.4 - 170)/543.1 = 0.0431$; $k_{sn} = 0.0431\times A_0^{0.5}
= \qty{43.1}{m}$.

\textbf{8.} Taksiran chi bertumpu pada ketinggian sepanjang lebih dari \qty{3}{km}
sungai, bukan pada dua belas gradien yang berisik.

\textbf{9.} $K = U/k_{sn} = 4.3\times10^{-5}/43.1 = \qty{1.00e-6}{yr^{-1}}$.

\textbf{10.} $10^{-6}\times\sqrt{3.87\times10^7}\times0.00691 =
\qty{4.30e-5}{m/yr}$, \qty{43}{m/Ma}: muara tetap mengimbangi
laju pengangkatan.

\textbf{11.} Di muara $10^{-6}\times6222 = \qty{6.2}{m/ka}$; di
air terjun, $d = \qty{7.40}{km}$, $A = \qty{1.70e7}{m^2}$ dan
\qty{4.1}{m/ka}.

\textbf{12.} $170 + 43\times(4.5)^{-0.5}(12000^{0.15} - 300^{0.15})/0.15 =
\qty{405}{m}$: relief sebesar \qty{235}{m}.

\textbf{13.} Untuk $n = 1$ persamaannya berupa gelombang kinematik: titik-titik di kedua
sisi undakan naik sebesar $U$ per satuan waktu, sehingga undakan mempertahankan tingginya,
yakni tinggi jatuhan sebesar \qty{25}{m}.

\textbf{14.} $\chi = (1000/\sqrt{4.5})(12000^{0.15} - 7400^{0.15})/0.15 =
\qty{899.4}{m}$.

\textbf{15.} $t = \chi/(KA_0^{0.5}) = 899.4/10^{-3} = \qty{899}{ka}$.

\textbf{16.} Pada $\chi = 899.4 + 300 = \qty{1199}{m}$:
$d^{0.15} = 12000^{0.15} - 1199\times0.15\times\sqrt{4.5}/1000$, $d =
\qty{6.25}{km}$, \qty{5.75}{km} di hulu muara.

\textbf{17.} Di pertemuan ($d = \qty{9.5}{km}$) $\chi = \qty{442.8}{m}$;
air terjun terletak di titik tempat $\chi$ anak sungai mencapai \qty{899.4}{m}:
$s^{0.15} = 5000^{0.15} - (899.4 - 442.8)\times0.15\times\sqrt{1.5}/1000$,
$s = \qty{4.27}{km}$, \qty{0.73}{km} di hulu pertemuan.

\textbf{18.} Pertemuan mereka berada di $\chi = 1954$ dan \qty{999}{m}, melampaui
\qty{899}{m}: \omterm{def:b3:quantitative-geomorphology:knickpoint}{titik tekuk} belum melewati keduanya, sehingga sungai utama di sana
masih pada ketinggian lamanya.

\textbf{19.} $E = 5\times10^{-3}\times0.0086 = \qty{4.3e-5}{m/yr}$,
\qty{43}{m/Ma}: bukit di atas air terjun terkikis pada laju pengangkatan, dalam
\omterm{def:b3:quantitative-geomorphology:steady-profile}{keadaan tunak}.

\textbf{20.} $150^2/(5\times10^{-3}) = \qty{4.5}{Ma}$, lima kali umur
air terjun.

\textbf{21.} Kaki lereng mereka diturunkan \qty{25}{m} bersama alur:
bagian bawahnya menjadi terlalu curam, mungkin meluncur, sementara puncaknya masih
terkikis pada laju lama. Gelombang kecuraman merambat naik pada setiap lereng; setelah beberapa
juta tahun lereng-lereng itu kembali ke parabola yang sama, \qty{25}{m} lebih rendah.

\textbf{22.} $t = \chi k_{sn}/(U A_0^{0.5})$ bervariasi sebagai $1/U$: dari
$899/1.2 = 749$ hingga $899/0.8 = \qty{1124}{ka}$.

\textbf{23.} Dengan $n = 2$ kecepatan rambat tumbuh bersama gradien: muka
yang curam akan merambat lebih cepat daripada ruas di sekitarnya dan tetap tajam, tetapi
kecepatannya dan umur yang dibaca darinya akan berbeda.

\textbf{24.} \omterm{def:b3:quantitative-geomorphology:knickpoint}{Dasar erosi} Sungai Varden turun
\textbf{\qty[text-series-to-math=true]{900}{ka}} yang lalu (\qty{899}{ka} menurut
perhitungan).
\end{solution}
